\documentclass[7pt,landscape]{article}
\usepackage[utf8]{inputenc}
\usepackage[margin=0.18in]{geometry}
\usepackage{multicol}
\usepackage{titlesec}
\usepackage{enumitem}
\usepackage{xcolor}

\definecolor{navy}{RGB}{0, 51, 102}
\definecolor{linuxc}{RGB}{160, 88, 24}
\definecolor{winc}{RGB}{45, 95, 140}
\setlength{\columnsep}{0.25in}
\setlist{nosep, leftmargin=10pt}
\titlespacing*{\section}{0pt}{1pt}{0.5pt}
\setlength{\parindent}{0pt}
\linespread{0.90}

\titleformat{\section}{\small\bfseries\color{navy}}{\thesection}{0.5em}{}[\titlerule]
\pagestyle{empty}

\begin{document}
\textbf{Role Reference: Engineer} \hfill \textit{Last Updated: Aug 2026}
\raggedcolumns
\begin{multicols*}{3}

\section{ROLE RESPONSIBILITIES}
Hardening at round start. Eviction and patching once a threat hunter hands off a confirmed foothold. Restoring a broken service without weakening it. Stripping the tooling red team's playbook depends on, ahead of any foothold.

\section{ORDER OF OPERATIONS}
\textbf{Round start:} Firewall, user cleanup, package updates, service config, red team tooling audit (below), in that order, per host. \\
\textbf{Nonessential services:} Start shut down. Restart only once validated. \\
\textbf{On handoff:} Confirm account, process chain, hole, and scope before starting. Ask if anything is missing. \\
\textbf{Evict:} Kill session, reset/remove account, remove persistence. \\
\textbf{Patch:} Fix the actual vulnerability, not just the artifact. Rotate anything possibly compromised. \\
\textbf{Protect the fix:} Harder to undo than it was to apply.

\section{DENY RED TEAM: SAFE TO STRIP}
Confirm nothing scored depends on it, then log the removal like any other change.
\begin{itemize}
    \item \textbf{Python 2/3, if unused:} Ansible's default connection plugin needs a Python interpreter on the managed node; most red team post-ex and staging scripts assume one too. \texttt{apt remove --purge python3} (or \texttt{chmod 000} the binary and \texttt{apt-mark hold} python3 to keep it pinned) breaks both.
    \item \textbf{Perl, Ruby, other unused interpreters:} same idea, lower prevalence. Rarely needed by a scored app, often needed by an older exploit or pivot script.
    \item \textbf{Compilers (gcc, cc, make, build-essential):} rarely needed at runtime. Blocks on-box compilation of a privilege-escalation exploit or custom implant.
    \item \textbf{netcat/ncat/socat, if unused:} the default pivot and reverse-shell relay. Check startup scripts and cron before removing.
    \item \textbf{SSH client (not server) on hosts that never SSH out:} removes a lateral-movement path, leaves inbound access alone.
    \item \textbf{curl/wget, if a service never calls out:} rarely safe to pull, but closes the fetch-and-execute pattern on a locked-down box.
\end{itemize}

\section{DENY RED TEAM: PROTOCOLS \& AGENTS}
\begin{itemize}
    \item \textbf{WinRM/PSRemoting:} transport for Ansible's \texttt{winrm} connection and common Empire/C2 lateral movement. Disable if unused: \texttt{Disable-PSRemoting -Force}, block 5985/5986 at the host firewall.
    \item \textbf{Admin shares (\texttt{C\$}, \texttt{ADMIN\$}):} what psexec/wmiexec ride in on. Scope to your management host with a firewall rule instead of a blanket disable if you use them yourself.
    \item \textbf{PowerShell v2:} installable alongside v5+, lacks v5's logging; attackers downgrade to it on purpose. \texttt{Disable-WindowsOptionalFeature -Online} \newline \texttt{-FeatureName MicrosoftWindowsPowerShellV2}.
    \item \textbf{Outbound access, generally:} cutting default outbound after hardening (Firewall Reference) kills fetch-and-execute, Ansible included.
    \item Needed by a scored service? Narrow it, don't cut it, and note why in the tracking template.
\end{itemize}

\section{EVICTION COMMANDS}
\begin{itemize}
    \item \textit{\textcolor{linuxc}{Linux:}} \texttt{kill -9 [PID]}. \textit{\textcolor{winc}{Windows:}} \texttt{Stop-Process -Id [PID] -Force}.
    \item Reset a stolen account's password immediately. Delete an attacker-created account.
    \item Remove persistence: cron (\texttt{crontab -e}, \texttt{/etc/cron.d}), scheduled tasks, unfamiliar \texttt{authorized\_keys}, any added service/startup entry.
    \item \texttt{userdel -r [user]} \textcolor{linuxc}{(Linux)} / \texttt{Remove-LocalUser} \textcolor{winc}{(Windows)}: remove an attacker-created account and its home directory.
    \item Log what you found before deleting it.
\end{itemize}

\section{PATCHING COMMANDS}
\begin{itemize}
    \item Update the vulnerable package: \texttt{apt update \&\& apt upgrade -y} \textcolor{linuxc}{(Debian)}, Windows Update / \texttt{DISM} \textcolor{winc}{(Windows)}.
    \item Close the specific exposed port or service.
    \item Fix the misconfiguration, not just its symptom.
    \item \texttt{debsums -c}: verify installed package files against known-good hashes \textcolor{linuxc}{(Debian)}, to confirm a patched binary is actually clean.
\end{itemize}

\section{SSH \& REMOTE ACCESS HARDENING}
\begin{itemize}
    \item \texttt{/etc/ssh/sshd\_config}: \texttt{PermitRootLogin no}, \texttt{PasswordAuthentication no} once keys are confirmed working, \texttt{AllowUsers [list]}.
    \item \texttt{systemctl restart sshd} after any config change; test a second session before closing the first.
    \item \texttt{fail2ban}: auto-ban repeated failed logins. \texttt{apt install fail2ban}, configure a jail for sshd.
    \item \textcolor{winc}{Windows}: NLA required for RDP, restrict Remote Desktop Users group membership.
\end{itemize}

\section{FIREWALL REFERENCE}
\begin{itemize}
    \item \texttt{ufw allow ssh}, \texttt{ufw allow [port]/tcp}, \texttt{ufw enable}. \textcolor{linuxc}{(Linux)}
    \item \texttt{firewall-cmd --permanent --add-service=http} then \texttt{--reload} \textcolor{linuxc}{(RHEL family)}.
    \item \textcolor{winc}{Windows}: \texttt{New-NetFirewallRule}, or \texttt{netsh advfirewall firewall add rule} for the legacy syntax.
    \item Disallow unnecessary outbound after hardening a host; most legitimate services do not need free outbound access.
\end{itemize}

\section{PROTECTING A FIX}
\begin{itemize}
    \item \texttt{chattr +i [file]}: makes a restored \textcolor{linuxc}{Linux} file immutable, even to root. Remove with \texttt{chattr -i} before intentional edits.
    \item \texttt{icacls} on \textcolor{winc}{Windows}: lock down write access tighter than default.
    \item Re-run the account/group baseline after any fix.
\end{itemize}

\section{BACKUP \& RESTORE}
\begin{itemize}
    \item \texttt{rsync -avz [src] [dest]}: fast, incremental-friendly copy. \textcolor{linuxc}{(Linux)}
    \item \texttt{robocopy [src] [dest] /MIR}: mirrors a directory tree. \textcolor{winc}{(Windows)}
    \item \texttt{vssadmin list shadows}: check for existing Volume Shadow Copies on \textcolor{winc}{Windows} before assuming nothing is recoverable.
    \item Always restore to a new location first if time allows, then swap in, rather than overwriting a possibly-still-needed original.
\end{itemize}

\section{VERIFY}
\begin{itemize}
    \item Re-run the file-integrity baseline against the same paths.
    \item Confirm the scored service is up and responding correctly.
    \item Watch for the same account, address, or pattern returning.
\end{itemize}

\section{PATCH MANAGEMENT}
\begin{itemize}
    \item \texttt{apt-mark hold [package]}: pin a package version after patching, so an unrelated update does not silently reintroduce a fixed issue.
    \item \texttt{yum versionlock [package]} (RHEL family): same idea.
    \item Confirm the patched version actually took: \texttt{dpkg -l [package]} or \texttt{rpm -q [package]} after any update.
    \item A restart may be required for some patches (kernel, some services) to actually apply; check before assuming a patch is live.
\end{itemize}

\section{WEB \& DATABASE HARDENING}
\begin{itemize}
    \item Disable directory listing: \texttt{autoindex off;} (Nginx) or \texttt{Options -Indexes} (Apache).
    \item Remove default/sample content shipped with the install; it is the first thing a scanner checks.
    \item Review database user privileges: an application account should never have superuser rights on the database it uses.
    \item Confirm the database is not listening on a public interface unless it genuinely needs to be: \texttt{bind-address} in MySQL, \texttt{listen\_addresses} in PostgreSQL.
\end{itemize}

\section{KERNEL \& OS-LEVEL HARDENING}
\begin{itemize}
    \item \texttt{sysctl net.ipv4.ip\_forward}: should be 0 unless the host is genuinely a router.
    \item \texttt{sysctl kernel.randomize\_va\_space}: confirm ASLR is enabled (value 2).
    \item Disable unused kernel modules that expand attack surface, if you know the host does not need them.
    \item \textcolor{winc}{Windows}: disable SMBv1 \newline (\texttt{Disable-WindowsOptionalFeature} \newline \texttt{-Online -FeatureName SMB1Protocol}) and LLMNR via Group Policy, unless still needed.
\end{itemize}

\section{SANITY CHECKS AFTER ANY CHANGE}
\begin{itemize}
    \item Re-test the service from a second session before closing the first, especially after an SSH or RDP config change.
    \item Confirm the firewall change did not also block something legitimate.
    \item Time-stamp the change in the tracking template immediately, not after moving to the next task.
\end{itemize}

\section{TASKS VS. SERVICES}
\begin{itemize}
    \item \texttt{systemctl list-units --type=service} \textcolor{linuxc}{(Linux)} or \texttt{Get-Service} \textcolor{winc}{(Windows)}: what is running now.
    \item Cross-reference against what the host should run.
    \item Keep a known-good config or backup of anything hardened, in a name and location that is not the obvious first guess.
\end{itemize}

\section{ONE-LINERS}
\begin{itemize}
    \item \texttt{find / -mmin -30 -type f 2>/dev/null}: files touched in the last 30 minutes, a fast recent-change sweep. \textcolor{linuxc}{(Linux)}
    \item \texttt{Get-ChildItem C:\textbackslash -Recurse | Where} \newline \texttt{LastWriteTime -gt (Get-Date).AddMinutes(-30)}: same idea. \textcolor{winc}{(Windows)}
    \item \texttt{systemctl show [service] -p ExecStart}: confirm exactly what a service launches, useful when a unit file looks tampered with.
    \item \texttt{netstat -tulnp | grep LISTEN}: quick view of everything actually listening right now.
\end{itemize}

\end{multicols*}
\end{document}